
\section{与拉格朗日子流形关联的滤过 $A_\infty$ 代数}
\label{sec:3-5}

本节说明，第 2 章的构造及其推广，会在 $M$ 的每个相对 spin 拉格朗日子流形
$L$ 上给出滤过 $A_\infty$ 代数的几何实现。本节的主要定理是定理 3.5.11。
我们再次强调，目前尚不对障碍类 $o_k(L)$ 施加
$[o_k(L)]=0$ 这一条件。本节取 $R=\mathbb Q$。

令 $S^k(L;\mathbb Q)$ 为由 $L$ 的余维 $k$ 光滑奇异单形生成的自由
$\mathbb Q$-模。令 $\mathcal X_1(L)$ 为光滑奇异单形的一个可数集，满足
\[
  \mathcal X_1(L)\supset\mathcal X_L,\qquad
  \mathcal X_1(L)\ne\mathcal X_L,
\]
其中 $\mathcal X_L$ 如定理 3.4.8 所示。以
$C^k(L;\mathbb Q)$ 表示 $S^k(L;\mathbb Q)$ 中由 $\mathcal X_1(L)$
生成的子模。（与上一节一样，我们采用“光滑对应”方法。为此必须使用横截性论证和
贝尔纲定理，这也正是要使用可数生成复形 $C^k(L;\mathbb Q)$ 的原因。
$C^k(L;\mathbb Q)$ 的构造见第 7.2 节。）与第 3.4 节一样，我们有时把
$\mathcal X_1(L)$ 的元素写成 $P=(|P|,f)$。定义
$C^\bullet(L;\Lambda_{0,\mathrm{nov}})$ 为
$C^\bullet(L;\mathbb Q)\otimes\Lambda_{0,\mathrm{nov}}$ 的完备化。
为简化记号，置
\[
  C^\bullet=C^\bullet(L;\Lambda_{0,\mathrm{nov}})
  :=\bigl(C^\bullet(L;\mathbb Q)
  \otimes\Lambda_{0,\mathrm{nov}}\bigr)^{\wedge}.
\]

回忆 $\Lambda_{0,\mathrm{nov}}$ 中元素 $T^\lambda e^\mu$ 的次数为 $2\mu$。
因此，$C^\bullet(L;\Lambda_{0,\mathrm{nov}})$ 中的次数，是
$C^\bullet(L;\mathbb Q)$ 中的次数与 $\Lambda_{0,\mathrm{nov}}$ 中系数次数
之和。利用 $\Lambda_{0,\mathrm{nov}}$ 上的滤过，可以唯一地在
$C^k(L;\Lambda_{0,\mathrm{nov}})$ 上定义满足下列条件的滤过：
\[
  C^k(L;\mathbb Q)\subset F^0C^k(L;\Lambda_{0,\mathrm{nov}})
\]
以及
\[
  C^k(L;\mathbb Q)\cap F^\lambda
  C^k(L;\Lambda_{0,\mathrm{nov}})=\{0\},
  \qquad \lambda>0.
\]

现在定义同态
\[
  \mathfrak m_k\colon
  B_k\bigl(C(L;\Lambda_{0,\mathrm{nov}})[1]\bigr)
  \longrightarrow C(L;\Lambda_{0,\mathrm{nov}})[1],
  \qquad k\geq0.
\]
为此，回忆 $\mathcal M_{k+1}(\beta)$ 是形如
$((\Sigma,\vec z),w)$ 的对子之集，其中
$(\Sigma,\vec z)\in\mathcal M_{k+1}$，而
\[
  w\colon(\Sigma,\partial\Sigma)\longrightarrow(M,L)
\]
是代表类 $\beta$ 的伪全纯映射。令
$\mathcal M_{k+1}^{\mathrm{main}}(\beta)$ 为
$\mathcal M_{k+1}(\beta)$ 中由 $(\Sigma,\vec z)$ 属于主分支的元素组成的子集
（见定义 2.1.20 和 2.1.27）。

给定 $P_i\in\mathcal X_1(L)$，考虑与上一节相同的纤维积
\begin{equation}
  \mathcal M_{k+1}^{\mathrm{main}}(\beta)
  {}_{(\operatorname{ev}_1,\ldots,\operatorname{ev}_k)}
  \!\times_{f_1\times\cdots\times f_k}
  (P_1\times\cdots\times P_k).
  \tag{3.5.1}\label{eq:3-5-1}
\end{equation}
令 $\vec P=(P_1,\ldots,P_k)$，其中
$P_i\in\mathcal X_1(L)$。纤维积 \eqref{eq:3-5-1} 是带仓西结构的空间；
以后把它记作
$\mathcal M_{k+1}^{\mathrm{main}}(\beta;\vec P)$。

\begin{proposition}{3.5.2}\label{prop:3-5-2}
设 $L$ 是相对 spin 拉格朗日子流形。则可以选取 $\mathcal X_1(L)$，
并在 $\mathcal M_{k+1}^{\mathrm{main}}(\beta;\vec P)$ 上选取一个多值截面系
$\mathfrak s_{\beta,\vec P}$，使得下列性质成立：
\begin{align}
  &\mathcal M_{k+1}^{\mathrm{main}}(\beta;\vec P)
  \text{ 的（虚）维数为}\notag\\[-0.3em]
  &\hspace{7em}
  n-\sum_i(g_i-1)+\mu_L(\beta)-2,
  \quad
  \begin{gathered}
    n=\dim L,\\[-0.2em]
    (P_i,f_i)\in C^{g_i}(L;\mathbb Q),
  \end{gathered}
  \tag{3.5.3.1}\label{eq:3-5-3-1}\\
  &\mathfrak s_{\beta,\vec P}\text{ 横截于 }0;
  \tag{3.5.3.2}\label{eq:3-5-3-2}\\
  &\mathfrak s_{\beta,\vec P}\text{ 在边界处相容};
  \tag{3.5.3.3}\label{eq:3-5-3-3}\\
  &\bigl(\mathfrak s_{\beta,\vec P}^{-1}(0),
  \operatorname{ev}_0\bigr)
  \text{ 是 }\mathcal X_1(L)\text{ 中元素的 }\mathbb Q\text{-线性组合}.
  \tag{3.5.3.4}\label{eq:3-5-3-4}
\end{align}
其中 $\mu_L(\beta)$ 是马斯洛夫指标。
\end{proposition}

更准确地说，(3.5.3.4) 的含义如下。由引理 A1.26，集合
$\mathfrak s_{\beta,\vec P}^{-1}(0)$ 具有光滑三角剖分。选定一个这样的剖分，
并固定其顶点的一个次序。这个剖分中每个维数为
\[
  d=\dim\mathfrak s_{\beta,\vec P}^{-1}(0)
\]
的单形 $\Delta_a^d$ 都带有重数
$\operatorname{mul}_{\Delta_a^d}$（见定义 A1.27）。把
$\operatorname{ev}_0$ 限制到 $\Delta_a^d$，便定义了光滑奇异单形
$(\Delta_a^d,\operatorname{ev}_0)$。(3.5.3.4) 的意思正是：每个
$(\Delta_a^d,\operatorname{ev}_0)$ 都属于 $\mathcal X_1(L)$。

虚基本链\liftIndex{I345}
$\bigl(\mathfrak s_{\beta,\vec P}^{-1}(0),\operatorname{ev}_0\bigr)$ 定义为
\[
  \bigl(\mathfrak s_{\beta,\vec P}^{-1}(0),\operatorname{ev}_0\bigr)
  =\sum_{a\in A}
  \pm\operatorname{mul}_{\Delta_a^d}
  (\Delta_a^d,\operatorname{ev}_0),
\]
其中 $\{\Delta_a^d\mid a\in A\}$ 是全部维数为
$d=\dim\mathfrak s_{\beta,\vec P}^{-1}(0)$ 的单形之集，正负号由定向决定
（见定义 A1.28）。

\begin{remark}{3.5.4}\label{rem:3-5-4}
应当强调，多值截面 $\mathfrak s_{\beta,\vec P}$ 的选取依赖于
$\beta$ 和 $\vec P$。我们必须选取满足某些相容条件 (3.5.3.3) 的
$\mathfrak s_{\beta,\vec P}$（见相容条件 7.2.38 和 7.2.44）。
还必须使 $\mathfrak s_{\beta,\vec P}^{-1}(0)$ 的三角剖分对于
$\beta,\vec P$ 的各种选择彼此相容。第 7.2 节将详细讨论这些问题。
\end{remark}

\begin{proof}
第 7 章将说明仓西结构的构造，并陈述命题 7.1.1、7.1.2 和 7.2.35；
这些命题说明如何选取上述 $C(L;\mathbb Q)$ 与多值截面 $\mathfrak s$。
关于 $\mathcal M_{k+1}^{\mathrm{main}}(\beta)$ 的定向以及仓西结构纤维积的
定向，见第 8 章的第 8.4、8.5 节。这里只证明 (3.5.3.1) 中的维数陈述。
直接计算得
\[
  \mu_L(\beta)+k+1-3+n-\sum_i g_i
  =\mu_L(\beta)+n-\sum_i(g_i-1)-2,
\]
即得所需结论。
\end{proof}

\begin{definition}{3.5.5}\label{def:3-5-5}
令多值截面 $\mathfrak s$ 和可数生成子模 $C(L;\mathbb Q)$ 如命题 3.5.2。
定义
\[
  \mathcal M_{k+1}^{\mathrm{main}}(\beta;\vec P)^{\mathfrak s}
  :=\bigl(\mathfrak s_{\beta,\vec P}^{-1}(0),
  \operatorname{ev}_0\bigr).
\]
（关于上述模空间的定向，见第 8 章的定义 8.4.1。）
\end{definition}

现在定义映射 $\mathfrak m_k$。回忆存在元素
$\beta_0=0\in G(L)$（$G(L)$ 的定义见第 3.1 节的 (3.1.8)），并且
\[
  \mu_L(\beta_0)=0,\qquad \omega(\beta_0)=0.
\]

\begin{definition}{3.5.6}\label{def:3-5-6}
\begin{enumerate}[label=\textup{(\arabic*)},leftmargin=2.8em]
  \item 对 $(P,f)\in C^k(L;\mathbb Q)$，定义
  \[
    \mathfrak m_{0,\beta}(1)=
    \begin{cases}
      \mathcal M_1(\beta)^{\mathfrak s},
        &\beta\ne\beta_0,\\
      0,&\beta=\beta_0,
    \end{cases}
  \]
  \[
    \mathfrak m_{1,\beta}(P,f)=
    \begin{cases}
      \mathcal M_2^{\mathrm{main}}(\beta;P)^{\mathfrak s},
        &\beta\ne\beta_0,\\
      (-1)^n\partial P,&\beta=\beta_0,
    \end{cases}
  \]
  以及
  \[
    \mathbf e=PD([L])\in C^0(L)=C[1]^{-1}(L),
  \]
  其中 $[L]$ 是 $L$ 的基本闭链，$PD$ 表示庞加莱对偶；
  $\mathfrak m_{1,\beta_0}$ 定义中的 $\partial$ 是通常的经典边界算子。

  \item 对每个 $k\geq2$ 及 $P_i\in C^{g_i}(L;\mathbb Q)$，定义
  \[
    \mathfrak m_{k,\beta}(P_1,\ldots,P_k)
    =\mathfrak m_{k,\beta}(P_1\otimes\cdots\otimes P_k)
    =\mathcal M_{k+1}^{\mathrm{main}}
      (\beta;\vec P)^{\mathfrak s}.
  \]

  \item 最后，对 $k\geq0$ 定义
  \begin{equation}
    \mathfrak m_k
    =\sum_{\beta\in G(L)}
    \mathfrak m_{k,\beta}\otimes
    T^{\omega(\beta)}e^{\mu_L(\beta)/2}.
    \tag{3.5.7}\label{eq:3-5-7}
  \end{equation}
\end{enumerate}
\end{definition}

\pagebreak[3]
\begin{remark}{3.5.8}\label{rem:3-5-8}
\begin{enumerate}[label=\textup{(\arabic*)},leftmargin=2.8em]
  \item 在上面对 $\mathfrak m_{1,\beta_0}$ 的定义中，我们把 $P$ 看作链，
  $\partial$ 是通常的边界算子。若把 $P$ 看作上链，则符号必须改为
  \[
    \mathfrak m_{1,\beta_0}(P)
    =(-1)^{n+\deg P+1}dP,
  \]
  其中 $\deg P$ 是 $P$ 移位前的次数。这里采用如下庞加莱对偶符号约定：
  对 $L$ 中的一条链 $S$，把庞加莱对偶 $PD(S)$ 看作流，则对任意
  $\alpha\in\Omega^{\dim S}(L)$ 都有
  \[
    \int_S\alpha|_S=\int_LPD(S)\wedge\alpha.
  \]
  注意，这一约定与第 8.6 节给出的法丛定向约定相容。在此约定下，
  \[
    PD(\partial S)=(-1)^{\deg S+1}d\bigl(PD(S)\bigr).
  \]

  \item 映射 $\mathfrak m_k$ 依赖于殆复结构 $J$ 以及仓西映射
  $\mathfrak s$ 的扰动（即障碍丛的一个多值截面；见第 A1 节）。
  第 4.6 节将详细研究这种依赖性。

  \item 因为 $\mathfrak m_{0,\beta_0}(1)=0$，条件 (3.2.12.7) 成立。

  \item 与 (3.2.13) 一样，可以把
  $(C(L;\Lambda_{0,\mathrm{nov}}),\mathfrak m)$ 的系数
  $\Lambda_{0,\mathrm{nov}}$ 约化为 $\mathbb Q$，得到
  $A_\infty$ 代数 $(\overline C,\overline{\mathfrak m})$；
  这是因为条件 (3.2.13.3) 是 (3.1.7.2) 的推论。又因
  $\mathfrak m_{0,\beta_0}=0$，此时 $\overline{\mathfrak m}_0=0$。
  还应注意，$(\overline C,\overline{\mathfrak m})$ 不是微分分次代数
  （D.G.A.），而是 $A_\infty$ 代数。考虑 $\beta=\beta_0(=0)$ 的情形。
  若 $P_1\ne P_2$，则作为集合，
  $\mathcal M_2^{\mathrm{main}}(\beta_0;P_1,P_2)$ 无非就是
  $P_1\cap P_2$。但关键在于，我们不能在链层面定义自交。类似地，也不能期待
  $\overline{\mathfrak m}_k=0$ 对 $k\geq3$ 成立，因为在对角情形没有横截性。
  第 7 章第 7.2 节将更详细地讨论这一问题。

  上一节为拉格朗日子流形 $L$ 构造了 $\mathbb Z$ 上的经典
  $A_\infty$ 代数
  $(\mathbb Z\mathcal X_L,\overline{\mathfrak m}^{\,L})$，
  并容许使用横截性论证。于是可以证明，
  $(C(L;\Lambda_{0,\mathrm{nov}}),\mathfrak m)$ 是
  $(\mathbb Q\mathcal X_L,\overline{\mathfrak m}^{\,L})$ 的
  $A_\infty$ 形变，其意义如定义 3.2.34（见下面的定理 3.5.11）。
  最后再说明一点符号。依照 (3.2.5)，还必须验证
  \begin{equation}
    \mathcal M_3^{\mathrm{main}}(\beta_0;P_1,P_2)
    =(-1)^{\deg P_1(\deg P_2+1)}P_1\cap P_2,
    \tag{3.5.9}\label{eq:3-5-9}
  \end{equation}
  其中 $\beta=\beta_0(=0)$。第 8 章的推论 8.6.4 将证明此式。
\end{enumerate}
\end{remark}

\begin{lemma}{3.5.10}\label{lem:3-5-10}
形式和 \eqref{eq:3-5-7} 收敛，并定义连续运算
\[
  \mathfrak m_k\colon
  B_k\bigl(C(L;\Lambda_{0,\mathrm{nov}})[1]\bigr)
  \longrightarrow C(L;\Lambda_{0,\mathrm{nov}})[1].
\]
\end{lemma}

\begin{proof}
由格罗莫夫紧性定理，$G(L)$ 满足条件 3.1.6。
\end{proof}

回忆我们置 $\deg\beta=\mu_L(\beta)$，并把 $(P,f)$ 的次数移位 $1$。
于是维数公式 (3.5.3.1) 蕴含：$\mathfrak m_k$ 在 $B(C[1])$ 上的次数为 $+1$。

下面是本节的主要定理。\liftIndexRangeStart{I104}

\begin{theorem}{3.5.11}\label{thm:3-5-11}
设 $L$ 是相对 spin 拉格朗日子流形，则有：
\begin{align}
  &(C(L;\Lambda_{0,\mathrm{nov}}),\mathfrak m)
  \text{ 是以 }\mathbf e\text{ 为同伦幺元的滤过 }A_\infty\text{ 代数};
  \tag{3.5.12.1}\label{eq:3-5-12-1}\\
  &(C(L;\Lambda_{0,\mathrm{nov}}),\mathfrak m)
  \text{ 是 }G(L)\text{-有隙的};
  \tag{3.5.12.2}\label{eq:3-5-12-2}\\
  &(C(L;\Lambda_{0,\mathrm{nov}}),\mathfrak m)
  \text{ 是定理 3.4.8 中经典 }A_\infty\text{ 代数}\notag\\[-0.3em]
  &\hspace{7em}
  (\mathbb Q\mathcal X_L,\overline{\mathfrak m}^{\,L})
  \text{ 的 }A_\infty\text{ 形变，意义如定义 3.2.34}.
  \tag{3.5.12.3}\label{eq:3-5-12-3}
\end{align}
\end{theorem}

(3.5.12.3) 的意思如下：像 (3.2.13) 那样把系数环
$\Lambda_{0,\mathrm{nov}}$ 约化为 $\mathbb Q$ 后，得到
$\mathbb Q$ 上满足 $\overline{\mathfrak m}_0=0$ 的 $A_\infty$ 代数
$(C(L;\mathbb Q),\overline{\mathfrak m})$；而
$\overline{\mathfrak m}$ 在
$\mathbb Q\mathcal X_L\subset C(L;\mathbb Q)$ 上的限制，恰好等于
第 3.4 节构造的 $\overline{\mathfrak m}^{\,L}$。

\begin{remark}{3.5.13}\label{rem:3-5-13}
如注 3.4.11 所说，我们不会直接用多值截面 $\mathfrak s$ 构造滤过
$A_\infty$ 结构。相反，固定任意能量 $E\geq0$ 和 $k$ 后，只使用能量不超过
$E$ 的伪全纯圆盘模空间构造一个 $A_{n,K}$ 结构。$A_{n,K}$ 结构的精确定义
见第 7.2.6 小节，构造细节见第 7.2.4--7.2.8 小节。
\end{remark}

\begin{proof}
性质 (3.2.12.1)--(3.2.12.5) 直接来自
$(C(L;\Lambda_{0,\mathrm{nov}}),\mathfrak m)$ 的构造。现在假定下面
(3.5.17) 中的符号规则，证明 $\{\mathfrak m_k\}_{k\geq0}$ 满足
$A_\infty$ 公式；符号将在第 8.5 节精确核对。

先证明 $\widehat d\circ\widehat d=0$。这归结为分析
$\mathcal M_{k+1}^{\mathrm{main}}(\beta;P_1,\ldots,P_k)$ 的边界
（见 (3.5.1)）。可以看出，该边界是
\[
  \sum_i
  \mathcal M_{k+1}^{\mathrm{main}}
  (\beta;P_1,\ldots,\partial P_i,\ldots,P_k)
\]
与下面图 3.5.1 所描述各项之和。

\begin{figure}[htbp]
  \centering
  \setlength{\unitlength}{0.9mm}
  \begin{picture}(100,102)
    \put(50,35){\circle{58}}
    \put(50,78){\circle{28}}
    \qbezier[14](39,91)(50,99)(62,91)
    \qbezier[16](22,18)(13,31)(20,47)
    \qbezier[16](79,20)(87,34)(80,49)
    \put(60.2,87.5){\line(1,-1){3.2}}
    \put(64,88){$P_i$}
    \put(36,78){\line(1,0){3.5}}
    \put(19,76){$P_{i+k_2-1}$}
    \put(69.5,55.5){\line(1,-1){3.2}}
    \put(74,55){$P_{i-1}$}
    \put(76,24){\line(-2,1){3.6}}
    \put(80,20){$P_1$}
    \put(50,6){\line(0,1){4}}
    \put(45,0){$\operatorname{ev}_0$}
    \put(31,13){\line(1,2){2.4}}
    \put(22,8){$P_k$}
    \put(22,43){\line(2,-1){3.6}}
    \put(4,44){$P_{i+k_2}$}
    \put(46,76){$\beta_2$}
    \put(46,34){$\beta_1$}
  \end{picture}
  \par\smallskip
  \renewcommand{\thefigure}{3.5.1}
  \refstepcounter{figure}
  \textbf{\figurename~\thefigure}\label{fig:3-5-1}
\end{figure}

另一方面，为证明 $\widehat d\circ\widehat d=0$，只需证明
\begin{equation}
\begin{aligned}
  &\sum_{\beta_1+\beta_2=\beta}
  \sum_{k_1+k_2=k+1}
  \sum_i
  (-1)^{\deg P_1+\cdots+\deg P_{i-1}+i-1}\\
  &\qquad{}\cdot
  \mathfrak m_{k_1,\beta_1}
  \bigl(P_1,\ldots,
  \mathfrak m_{k_2,\beta_2}
  (P_i,\ldots,P_{i+k_2-1}),\ldots,P_k\bigr)=0.
\end{aligned}
\tag{3.5.14}\label{eq:3-5-14}
\end{equation}
（见注 3.2.21 (1)。）为简化记号，这里以 $P_i$ 代替 $(P_i,f_i)$。
把 \eqref{eq:3-5-14} 的左端依照
$\beta_1=0$ 且 $k_1=1$、$\beta_2=0$ 且 $k_2=1$，以及其余情形分成三项。
\liftIndexRangeEnd{I104}
于是可以把 \eqref{eq:3-5-14} 的左端改写成
\begin{equation}
\begin{aligned}
  &\mathfrak m_{1,0}\mathfrak m_{k,\beta}(P_1,\ldots,P_k)\\
  &\quad{}+\sum_i
  (-1)^{\deg P_1+\cdots+\deg P_{i-1}+i-1}
  \mathfrak m_{k,\beta}
  (P_1,\ldots,\mathfrak m_{1,0}(P_i),\ldots,P_k)\\
  &\quad{}+
  \sum_{\substack{
    \beta_1+\beta_2=\beta,\ k_1+k_2=k+1;\\
    \beta_1\ne0\ \mathrm{or}\ k_1\ne1,\\
    \beta_2\ne0\ \mathrm{or}\ k_2\ne1}}
  \sum_i
  (-1)^{\deg P_1+\cdots+\deg P_{i-1}+i-1}\\
  &\hspace{7em}{}\cdot
  \mathfrak m_{k_1,\beta_1}
  \bigl(P_1,\ldots,
  \mathfrak m_{k_2,\beta_2}
  (P_i,\ldots,P_{i+k_2-1}),\ldots,P_k\bigr).
\end{aligned}
\tag{3.5.15}\label{eq:3-5-15}
\end{equation}
由定义 3.5.6，
$\mathfrak m_{1,0}=(-1)^n\partial$，其中 $\partial$ 是经典边界映射。
因此，\eqref{eq:3-5-15} 的第一项无非是
\begin{equation}
  (-1)^n\partial
  \mathcal M_{k+1}^{\mathrm{main}}
  (\beta;P_1,\ldots,P_k),
  \tag{3.5.16.0}\label{eq:3-5-16-0}
\end{equation}
第二项是下式对 $i$ 的求和：
\begin{equation}
  (-1)^{n+\sum_{j=1}^{i-1}(\deg P_j+1)}
  \mathcal M_{k+1}^{\mathrm{main}}
  (\beta;P_1,\ldots,\partial P_i,\ldots,P_k).
  \tag{3.5.16.1}\label{eq:3-5-16-1}
\end{equation}
第三项在几何上对应于图 3.5.1 所描述的模空间，也就是下式之和：
\begin{equation}
\begin{aligned}
  &(-1)^{\sum_{j=1}^{i-1}(\deg P_j+1)}
  \mathcal M_{k-k_2+2}^{\mathrm{main}}
  \bigl(\beta_1;P_1,\ldots,\\
  &\hspace{7em}
  \mathcal M_{k_2+1}^{\mathrm{main}}
  (\beta_2;P_i,\ldots,P_{i+k_2-1}),\ldots,P_k\bigr).
\end{aligned}
\tag{3.5.16.2}\label{eq:3-5-16-2}
\end{equation}

此外，关于这些空间的定向，可以证明
\begin{equation}
\begin{aligned}
  &(-1)^n\mathcal M_{k+1}^{\mathrm{main}}
  (\beta;P_1,\ldots,\partial P_i,\ldots,P_k)\\
  &\quad\subseteq
  (-1)^{n+1+\sum_{j=1}^{i-1}(\deg P_j+1)}
  \partial\mathcal M_{k+1}^{\mathrm{main}}
  (\beta;P_1,\ldots,P_k)
\end{aligned}
\tag{3.5.17.1}\label{eq:3-5-17-1}
\end{equation}
以及
\begin{equation}
\begin{aligned}
  &\mathcal M_{k-k_2+2}^{\mathrm{main}}
  \bigl(\beta_1;P_1,\ldots,
  \mathcal M_{k_2+1}^{\mathrm{main}}
  (\beta_2;P_i,\ldots,P_{i+k_2-1}),\ldots,P_k\bigr)\\
  &\quad\subseteq
  (-1)^{n+1+\sum_{j=1}^{i-1}(\deg P_j+1)}
  \partial\mathcal M_{k+1}^{\mathrm{main}}
  (\beta;P_1,\ldots,P_k).
\end{aligned}
\tag{3.5.17.2}\label{eq:3-5-17-2}
\end{equation}
(3.5.17.1) 和 (3.5.17.2) 将在第 8 章的命题 8.5.1 中证明。
由此可见，(3.5.16.0) 与 (3.5.16.1)、(3.5.16.2) 之和相互抵消，
所以 \eqref{eq:3-5-15} 等于零。这就证明了
$\widehat d\circ\widehat d=0$。(3.5.12.2) 由构造立即可得。

为了构造同伦幺元，还必须研究横截性；证明推迟到第 7.3 节
（符号则见第 8.10 节）。除 (3.5.9) 的符号要在推论 8.6.4 中核对外，
(3.5.12.3) 直接来自第 7.2.4--7.2.7 小节将要给出的构造。
\end{proof}
